\scanpage{264}
\mspage{433}

\[
(ad-bc)D=(c^2+cd+d^2) \tag{23}
\]
donc $D$ \uline{est inversible \unsure{ssi}} $c^2+\unsure{cd|d}+d^2=\mathcal{N}_{A_\rho/A}(c\rho+d)$ \uline{l'est i.e.} \unsure{ssi} $c\rho+d$ est inversible, et alors
\[
ad-bc=(c^2+cd+d^2)/D
\]
\margin{car $c^2+cd+d^2\equiv1\ (3)$}
ce qui montre que le signe de $\det\alpha=ad-bc$ est alors égal à la valeur de $D$ mod $3$.

(Mais est-ce que $D$ peut être égal à $0$ mod $3$ ? alors mod $3$ la matrice $\xi$ \deleted{serait} dépendrait des seuls paramètres $A,B$
\[
\xi=\begin{pmatrix}A&B\\1+B&0\end{pmatrix}
\]
avec $-B(1+B)=1$ \inserted{i.e. $B^2+B+1=0$} \deleted{i.e. $B=0$ ou $B=-1$} \inserted{donc $B^3=1$ donc (car $3$) $B=1$}
donc on aura $\xi=\begin{pmatrix}A&1\\-1&0\end{pmatrix}$, \deleted{\illegible{4}}
\[
\xi\rho=\begin{pmatrix}-1&A+1\\0&-1\end{pmatrix},
\]
cette matrice est régulière \unsure{ssi} $A\ne-1$, OK.\unsure{)|\}} Donc a priori on ne peut exclure le cas où $D$ \unsure{peut} prendre une valeur nulle mod $3$\ldots)

Lorsqu'on prend la section $\sigma$ de $\unsure{\mathrm{Sl}(2)_{\mathbb{Z}}|\mathrm{Gl}(2)_{\mathbb{Z}}}$, donnée par
\[
\sigma=\sigma_\infty=\begin{pmatrix}0&-1\\1&0\end{pmatrix}
\qquad \sigma_\infty^4=1,\quad \sigma_\infty^2=-1 \tag{24}
\]
on a encore (sauf erreur)
\margin{\unsure{cas|car} $A=\mathbb{Z}$.}
\[
\mathrm{Pic}(A_\sigma)=\unsure{0|1}. \tag{25}
\]

\uline{La} condition numérique sur
\[
\eta=\begin{pmatrix}R&S\\T&U\end{pmatrix} \tag{26}
\]
pour \deleted{satisfaire} être \unsure{localmt} dans l'image

\scanpage{265}

de $F_\sigma$ (en plus de la condition de \uline{régularité} sur $\eta\sigma$)
\[
\eta=\beta\sigma(\beta)^{-1}=[\beta,\sigma] \tag{27}
\]
c'est qu'on ait
\[
\det\eta=1,\quad \mathrm{Tr}(\eta\sigma)=\mathrm{Tr}\sigma \tag{28}
\]
or
\[
\eta\sigma=\begin{pmatrix}R&S\\T&U\end{pmatrix}
\begin{pmatrix}0&-1\\1&0\end{pmatrix}
=\begin{pmatrix}S&-R\\U&-T\end{pmatrix}
\]
et (28) devient
\[
RU-ST=1,\quad S-T=0\quad\text{i.e. }S=T \tag{29}
\]
i.e. $\eta$ doit être \uline{symétrique}. \deleted{\illegible{1}} Si cette condition est satisfaite, il s'ensuit \unsure{par} (25) --- quand la base est $\mathbb{Z}$ --- que l'on peut trouver une solution $\beta\in\mathrm{Gl}(2,\mathbb{Z})$ de (27), l'obstruction à trouver une solution dans $\mathrm{Sl}(2,\mathbb{Z})$ est encore un signe, savoir un élément dans
\[
\mathbb{Z}^*/\mathcal{N}(A_\sigma^*)\simeq\pm1
\]
--- comme on a, pour
\[
\beta=\begin{pmatrix}r&s\\t&u\end{pmatrix}. \tag{30}
\]
\deleted{\unsure{la relation}} satisfaisant (27), la relation\margin{\unsure{$\varepsilon$}}
\[
(\unsure{rt-su})U=t^2+u^2=\mathcal{N}(t\sigma+u) \tag{31}
\]

\scanpage{266}
\mspage{434}

on voit comme tantôt que $U$ est inv. \unsure{ssi}\margin{(base quelc.)} $t\sigma+u$ l'est i.e. $t^2+u^2$ l'est, et que dans ce cas, l'obstruction \deleted{\unsure{de}} (base $\mathbb{Z}$, ou $\mathbb{Z}_\ell$) est le reste de $U$ mod $4$ (qui est égal à $\pm1$)\ldots
\margin{car $t^2+u^2\equiv1\ (4)$}

Résumons pour mémoire

\uline{Proposition.} Soit $A=\mathbb{Z}$, ou $\mathbb{Z}_\ell$, ou $\widehat{\mathbb{Z}}$, et considérons les \inserted{deux} équations \deleted{dans}
\[
\left\{\begin{aligned}
\relax
[\alpha,\rho]&\overset{\unsure{\mathrm{déf}}}{=}\alpha\rho(\alpha)^{-1}=\xi\\
[\beta,\sigma]&=\beta\sigma(\beta)^{-1}=\eta
\end{aligned}\right.
\]
où $\rho=\begin{pmatrix}0&1\\-1&1\end{pmatrix}$, $\sigma=\begin{pmatrix}0&-1\\1&0\end{pmatrix}$,
$\xi$ et $\eta\in\mathrm{Gl}(2,A)$ donnés, d'inconnues $\alpha$ resp. $\beta$ dans $\mathrm{Gl}(2,A)$. On suppose que
\[
\xi=\begin{pmatrix}A&B\\C&D\end{pmatrix}\quad\text{resp. }\eta=\begin{pmatrix}R&S\\T&U\end{pmatrix}
\]
satisfont les conditions \inserted{néc.} numériques pour l'existence d'une solution
\[
AD-BC=1\qquad B+D-C=1 \tag{16}
\]
\[
RU-ST=1\qquad S-T=0\quad\text{i.e. }S=T. \tag{29}
\]

Alors :

1) Pour qu'il existe une solution \inserted{$\alpha$ (resp. $\beta$)} dans $\mathrm{Gl}(2,A)$, il \unsure{faut|f.} et suffit que $\xi\rho=\begin{pmatrix}-B&A+B\\-D&C+D\end{pmatrix}$ soit régulière, i.e. non scalaire sur chaque fibre (resp. que $\eta\sigma=\begin{pmatrix}S&-R\\U&-T=-S\end{pmatrix}$ le soit) --- \deleted{et pour}

\scanpage{267}

\deleted{\unsure{ceci il suffit}} i.e. que sur aucune fibre, $\xi$ (resp. $\eta$) ne soit un multiple scalaire de $\rho^{-1}$ (resp. $\sigma^{-1}$) --- et pour ceci il suffit que $D$ (resp. $U$) soit inversible (car le \deleted{\unsure{coeff}} \unsure{correspondant} \unsure{dans|de}
\[
\rho^{-1}=\begin{pmatrix}1&-1\\1&0\end{pmatrix}\quad\text{ou}\quad
\sigma^{-1}=\begin{pmatrix}0&1\\-1&0\end{pmatrix}
\]
est zéro).

2) Lorsque $D$ (resp. $U$) est inversible, il y a tjrs une solution $\alpha$ resp. $\beta$, avec $\det\alpha\in\pm1$ (resp. $\det\beta\in\{\pm1\}$).

3) Pour qu'il y ait une solution dans $\mathrm{Sl}(2,A)$, il f. et s. qu'on ait $D\equiv1\ (3)$ (resp. $U\equiv1\ (4)$) --- condition qui est automatique si $A=\mathbb{Z}_\ell$ avec $\ell\ne3$ (resp. avec $\ell\ne2$).

II) \uline{Étude \inserted{dans $\mathrm{Gl}(2)$ \deleted{\unsure{$\mathbb{Z}$}}} de l'équation}
\[
[\alpha,\rho]=[\beta,\sigma]\unsure{\varepsilon_1}(\deleted{\illegible{1}}\mu-1). \tag{32}
\]
\margin{Étude \unsure{d'abord} de l'équation en $\xi,\eta$.}

On suppose tjrs que $G$ soit une forme de $\mathrm{Gl}(2)$, associé à une alg. d'Azumaya $\mathcal M$ sur la base $S=\mathrm{Spec}\,A$, et qu'on a \uline{deux} sections $\rho,\sigma$ \uline{régulières} de $G$ i.e. deux sections inv. de $\mathcal M$. On se donne de plus un ``sous-groupe à $1$ paramètre'' (additif) de $G$, i.e. un homo
\[
\varepsilon_1:\mathbf G_a\longrightarrow \mathrm{S}G,\qquad \lambda\longmapsto\varepsilon_1(\lambda). \tag{33}
\]

\scanpage{268}
\mspage{435}

Le cas qui nous intéresse plus particulièrement est celui où $G=\mathrm{Gl}(2)_{\mathbb{Z}}$, $\rho$ et $\sigma$ comme dessus, et
\[
\varepsilon_1(\lambda)=\varepsilon_1^\lambda=\begin{pmatrix}1&0\\\lambda&1\end{pmatrix}. \tag{34}
\]
On notera qu'ici
\[
\varepsilon_1=\varepsilon_1(1)=\begin{pmatrix}1&0\\1&1\end{pmatrix}=\deleted{\unsure{\ell_1}}\rho\sigma \tag{35}
\]
Nous considérons (32) comme des équations reliant les quantités $\alpha,\beta,\mu$ --- $\alpha,\beta$ sections \inserted{locales variables} de $G$, $\mu$ section de $\unsure{\mathbb E^1|\mathbf E^1}$ --- les hypothèses que nous allons faire \unsure{impliquent|impliqueront} d'ailleurs que $\mu$ sera en fait une section de $\mathbf G_m$ (s'il satisfait avec $\alpha,\beta$ à la relation (32)). \deleted{\illegible{1}}

On commencera par étudier les syst. \deleted{de solutions} $(\xi,\eta,\unsure{\mu|\beta})$ satisfaisant
\[
\left\{\begin{aligned}
\xi&=\eta\varepsilon_1(\mu-1) && (33)\\
\det\xi&=\det\eta=1 && (34)\\
\mathrm{Tr}\xi\rho&=\underbrace{\mathrm{Tr}\rho}_{\inserted{\substack{\Vert\\\unsure{i}}}}\ \unsure{\text{et}}\quad
\mathrm{Tr}\eta\sigma=\underbrace{\mathrm{Tr}\sigma}_{\inserted{\substack{\Vert\\j}}} && (35)
\end{aligned}\right.
\]
\margin{\uline{NB} On aura $\det\varepsilon_1(\mu-1)=1$ donc $\det\xi=\det\eta$, \unsure{par la première relation}.}

On peut éliminer $\eta$ de ces équations, grâce à (33), on trouve

\scanpage{269}

\[
(36)\qquad
\det\xi=1,\qquad
\operatorname{Tr}(\xi\inserted{-1})\rho
=\deleted{\unsure{i}}\inserted{0},\qquad
\unsure{\operatorname{Tr}(\xi\varepsilon_1(1-\mu)\deleted{\sigma}\inserted{-1})\inserted{\sigma}}
\deleted{=\unsure{j}}=0.
\]

\deleted{où on a \unsure{posé}, pour \unsure{simplifier}
\[
i=\operatorname{Tr}\rho\qquad j=\operatorname{Tr}\sigma
\]}

La troisième relation (36), en \unsure{développant}
\unsure{p.r. aux} coefficients de $\xi$, \deleted{et} peut s'écrire
\[
(37)\qquad P(\xi)\mu\unsure{-|+}Q(\xi)=0
\]
où les coefficients $P,Q$ sont des fonctions linéaires affines p.r. à ceux de
$\xi$, \deleted{et} \unsure{donnent} lieu à l'identité (en $\xi,\mu$)
\[
(38)\qquad
\unsure{\bigl(\operatorname{Tr}(\xi\varepsilon_1(1-\mu)\deleted{\sigma}\inserted{-1})\inserted{\sigma}\deleted{-j}\bigr)}
=P(\xi)\mu-Q(\xi)\deleted{\unsure{=0}}
\]
\deleted{\unsure{$\operatorname{Tr}(\sigma\xi\varepsilon_1(1-\mu))$}}

Le terme constant $-Q(\xi)$ étant \deleted{\unsure{égal à}}
\inserted{identiquement} \inserted{égal à} la valeur du premier
\unsure{membre} de (38) pour $\mu=0$, i.e.
\[
\begin{aligned}
(39)\qquad -Q(\xi)
&=\operatorname{Tr}((\xi\varepsilon_1-1)\sigma).\\
&=\operatorname{Tr}\xi(\varepsilon_1\sigma)-\operatorname{Tr}\sigma
\end{aligned}
\]

On suppose qu'on a
\[
(40)\qquad \varepsilon_1\sigma=-\lambda_0\rho
\qquad
\begin{gathered}
\lambda\text{ section centrale de }G\\
\text{i.e. }\lambda\text{ scalaire inversible}
\end{gathered}
\]
de sorte que
\[
(40\text{ bis})\qquad \boxed{\varepsilon_1=-\lambda_0\rho\sigma^{-1}}
\]
\margin{\uline{NB} Si $\det\rho=\det\sigma=1$,
alors on doit avoir $\lambda_0^2=1$
donc si $A$ \unsure{intègre}, $\lambda_0\in\{\pm1\}$.}
et (39) devient, en vertu \deleted{de \unsure{l'éq}}
\[
\deleted{\unsure{-}Q(\xi)=}\qquad
\operatorname{Tr}\xi(\varepsilon_1\sigma)
=\unsure{-\lambda_0|\lambda_0}\operatorname{Tr}\xi\rho
\overset{\inserted{\text{par }(36)}}{=}
-\lambda_0\operatorname{Tr}\rho,
\quad\text{donc}
\]

\scanpage{270}
\mspage{436}

on trouve
\[
\deleted{(\unsure{42|37})}\ (41)\qquad
Q(\xi)=\underbrace{\lambda_0\operatorname{Tr}\rho+\operatorname{Tr}\sigma}_{c}
\qquad
(\text{constante, indépendante de }\xi)
\]
et on suppose
\[
(42)\qquad
c=\lambda_0\operatorname{Tr}\rho+\operatorname{Tr}\sigma
\quad
\deleted{\text{\unsure{entier|sect.} \unsure{élément}}}
\inserted{\text{section}}\text{ \uline{inversible} de }\mathcal O_S.
\]

(Dans le cas qui nous intéresse, on aura
\[
\sigma^{-1}=-\sigma,\quad
\varepsilon_1=\rho\sigma=-\rho\sigma^{-1},\quad
\text{i.e. }\lambda_0=1
\]
\[
\deleted{\unsure{\operatorname{Tr}}}\quad
c=\lambda_0\underbrace{\operatorname{Tr}\rho}_{1}
+\underbrace{\operatorname{Tr}\sigma}_{0}=1
\qquad\text{OK.}
\]
)

Ceci posé, la \inserted{\uline{troisième}} équation (36) devient
\unsure{(compte tenu des deux premières)}
\[
(43)\qquad P(\xi)\mu=c
\]
ce qui signifie aussi que $P(\xi)$ \uline{est inversible}, et
\[
(44)\qquad \mu=c/P(\xi).
\]

Ainsi, $\mu$ est déterminé de façon unique par $\xi$, par l'équation (44) ---
et $\mu$ sera nécessairement inversible.

En résumé :

\uline{Proposition} Si on suppose qu'on a \deleted{(40 bis)}
une expression (40 bis) pour l'élément unipotent
$\varepsilon_1=\varepsilon_1(1)$, et que l'on ait (42)
($\lambda_0\operatorname{Tr}\rho+\operatorname{Tr}\sigma$ inversible), alors

\scanpage{271}

l'ens des solutions \deleted{\unsure{des|en}} $(\xi,\eta,\mu)$ des
équations simultanées (33)(34)(35)
est en corr. $1$-$1$ avec l'ens des $\xi$
\unsure{sections|section} de $G=G(\mathcal{M})$
satisfaisant les conditions
\[
(45)\qquad
\left\{
\begin{gathered}
\det\xi=1,\quad \operatorname{Tr}\xi\rho=\operatorname{Tr}\rho,\quad\text{et}\\
P(\xi)\text{ section inversible de }\mathcal O_S.
\end{gathered}
\right.
\]

En termes de $\xi$, la solution corr. du \unsure{syst.} (33 : 35)
est donnée par
\[
(\xi,\ \eta=\xi\varepsilon_1(1-\mu),\ \mu=c/P(\xi))
\]
--- et $\mu$ (le “multiplicateur”) est nécessairement inversible.

\uline{NB} Si on avait résolu en $\eta$, on aurait
trouvé une condition de la forme
\[
P'(\eta)\mu-Q'(\eta)=0
\]
avec
\[
Q'(\eta)=\operatorname{Tr}\eta\varepsilon_0^{-1}\rho-\operatorname{Tr}\rho
\]
et la condition naturelle serait
\[
\varepsilon_0^{-1}\rho=-\lambda'_0\sigma
\qquad(\lambda'_0\text{ scalaire inv}),
\]
qui est équivalente à (40 bis) avec $\lambda'_0=1/\lambda_0$,
et implique que
\[
Q'(\eta)=c^{\mathrm{te}}
=\lambda'_0\operatorname{Tr}\sigma-\operatorname{Tr}\rho
=\deleted{\unsure{\lambda'_0}}c',
\]

\scanpage{272}
\mspage{437}

on voit tout de suite que
\[
c'=\lambda'_0c
\]
donc la condition $c$ inv. équivaut à $c'$ inv.,
et on aura alors le système d'équations en $\eta$
\[
\det\eta=1,\quad
\operatorname{Tr}\eta\sigma=\operatorname{Tr}\sigma,\quad
P'(\eta)\mu=c'\quad\text{i.e. }
\left\{
\begin{gathered}
P'(\eta)\text{ inv.}\\
\mu=c'/P'(\eta)
\end{gathered}
\right.
.
\]

On voit donc que les hypothèses faites sont en fait symétriques en
$\xi,\eta$, et permettent un traitement symétrique de la résolution,
soit via $\xi$, soit via $\eta$.

Détermination de $P(\xi)$. On aura
\[
\varepsilon_1(\lambda)=1+N_1\lambda=(1+N_1)^\lambda
\]
$N_1$ une matrice de carré nul
\[
N_1^2=1,
\]
et
\[
\operatorname{Tr}(\xi\varepsilon_1(1-\mu)\sigma)-\operatorname{Tr}\sigma
=\operatorname{Tr}\xi\sigma-\operatorname{Tr}\sigma
+\unsure{\lambda}(1-\mu)\operatorname{Tr}\xi N_1\sigma
\]
donc
\[
(47)\qquad
P(\xi)=\deleted{\unsure{-}}\operatorname{Tr}\xi
\underbrace{(-N_1\sigma)}_{\inserted{\substack{\Vert\\N_1\rho=\rho+\sigma}}}
\]
D'ailleurs \deleted{la formule} on aura
\[
(48)\qquad \varepsilon_1=1+N_1
\]
et la formule (40 bis) s'écrit
\[
\deleted{\sigma+N_1\sigma=-\lambda_0\rho}
\]
\[
(49)\qquad N_1=-(\lambda_0\rho\sigma^{-1}+1)
\]
\margin{\unsure{i.e.}}
\[
\deleted{N_1\sigma=-\lambda_0\rho-\sigma}
\]
\deleted{soit} \unsure{ou} encore

\scanpage{273}

\[
N_1\sigma=-(\lambda_0\rho+\sigma)
\]
donc on aura
\[
(50)\qquad
P(\xi)=\operatorname{Tr}\deleted{\unsure{\xi}}\xi(\lambda_0\rho+\sigma)
\]

\uline{Corollaire de la Proposition} Le schéma des solutions
\inserted{des équ.} simultanées (33) : (35)
\inserted{en $\xi,\eta,\mu$} est un schéma relatif lisse,
de dim relative $2$ --- car l'équation $\operatorname{Tr}(\xi\rho)=c^{\mathrm{te}}$
définit toujours un sous-schéma lisse de dim $1$
de \deleted{\unsure{$\mathrm{Sl}(2)$|SG}} $SG(\mathcal{M})$
(dès que $\rho$ est une section qui n'est nulle sur aucune fibre),
--- dont il faut prendre ensuite l'ouvert défini par
“$P(\xi)$ inversible”.

Dans le cas qui nous intéresse,
\[
\rho=\begin{pmatrix}0&1\\-1&1\end{pmatrix},\qquad
\sigma=\begin{pmatrix}0&-1\\1&0\end{pmatrix},
\]
\deleted{la condition en $\xi$ \illegible{1} $P(\xi)$ inv. s'écrit}

\unsure{posant|pour}
\[
\xi=\begin{pmatrix}A&B\\C&D\end{pmatrix},\qquad
\eta=\begin{pmatrix}R&S\\T&U\end{pmatrix}
\]
on aura
\[
P(\xi)=D,\qquad P'(\eta)=U
\]
\[
(51)\qquad
\unsure{\eta}\varepsilon_1(\mu-1)
=\begin{pmatrix}
R+(\mu-1)S,&S\\
T+(\mu-1)U&U
\end{pmatrix}
\]

\scanpage{274}
\mspage{438}

et l'équation
\[
\xi=\eta\varepsilon_1(\mu-1)
\]
\deleted{\unsure{équivaut à}} implique
\[
(52)\qquad B=S,\quad D=U,
\]
et en fait elle \unsure{équivaut} à l'une de ces deux équations et à
\[
(53)\qquad \mu D=1\qquad (\unsure{\Longleftrightarrow|\Longrightarrow}\ \mu U=1)
\]
quand on suppose déjà que $\xi,\eta$ \unsure{satisfont} à l'équation des \deleted{traces et des} déterminants \inserted{et des traces}
\[
(54)\quad
\left\{\begin{array}{ll}
AD-BC=1 & RU-ST=1\\
\unsure{C|A}+D-B=1 & T-S=1
\end{array}\right.
\]

\unsure{Le} \deleted{\illegible{1}} schéma des solutions des équations simultanées $(33)$ \unsure{à|:} $(35)$ est alors celui des $(D,B)$ avec $D$ inv., ou encore des $(S,U)$ avec $S$ inv. On peut l'identifier au schéma des matrices
\[
(55)\qquad
\begin{pmatrix}\mu&M\\0&1\end{pmatrix}
=\begin{pmatrix}1/D&B/D\\0&1\end{pmatrix}
\]
(cf p.~410 \illegible{1}), qui \unsure{normalisent} le sous-groupe à $1$ \unsure{paramètre}
\[
\deleted{\unsure{\rho}\,\unsure{\varepsilon}}\quad
\varepsilon_0(\unsure{\lambda})=\deleted{\illegible{1}}
\begin{pmatrix}1&-\lambda\\0&1\end{pmatrix}
=\rho^{-1}\varepsilon_1(\lambda)\rho.
\]


\scanpage{275}

Revenons au cas général, où on suppose néanmoins que $\varepsilon_1$ donné par $(40\ \mathrm{bis})$, \deleted{\unsure{alors}} posons
\[
(56)\qquad \varepsilon_0=\rho^{-1}\varepsilon_1\rho
=-\lambda_0\sigma^{-1}\rho=1+\unsure{N_0},\quad\text{avec}
\]
\[
(57)\qquad N_0=\rho^{-1}N_1\rho=-(\lambda_0\sigma^{-1}\rho+1),
\]
de sorte que $N_0^2=0$, et
\[
(58)\qquad \varepsilon_0\quad\lambda\longmapsto
\varepsilon_0(\lambda)\overset{\unsure{\mathrm{dfn}}}{=}1+\lambda N_0
\]
est un sous-groupe à $1$ \unsure{paramètre} de $SG(\mathcal M)$, soit $L_0$, dont le \unsure{normalisateur} dans $G(\mathcal M)$ sera noté $B_0$. On a deux \unsure{homs}

$(59)$
\diagram{\begin{tikzcd}[ampersand replacement=\&]
B_0 \arrow[r, shift left=1ex, "\chi(\underline u)=\operatorname{dét}(\underline u)"]
\arrow[r, shift right=1ex, "\chi'(\underline u)\ \text{donné}"'] \& \unsure{\mathbf G_m|G_m}
\end{tikzcd}}
par
\[
\underline u(\varepsilon_0(\lambda))=\varepsilon_0(\chi'(\underline u)\lambda)
\]
ou encore (en termes linéaires dans $\mathcal M$)
\[
(60)\qquad \underline uN_0\underline u^{-1}=\chi'(\underline u)N_0
\]
\margin{Je souligne $\underline u$, pour ne pas confondre avec un coefficient de $\unsure{\eta|\beta}=\begin{pmatrix}r&s\\t&u\end{pmatrix}$ !}
Et on \unsure{considère} le sous-groupe
\[
(61)\quad\left\{\begin{array}{l}
\widetilde B_0^{\unsure{*}}\subset B_0\quad\text{défini par }\chi(\underline u)=\chi'(\underline u)\\
\text{i.e. par }\deleted{\unsure{\det}}\\
\underline uN_0\underline u^{-1}=\det(\underline u)N_0
\end{array}\right.
\]

donc on a
\[
L_0\subset\widetilde B_0^*.
\]
et $\overline{\mathrm m}$ on a une suite exacte de sch. en gr.
\diagram{\begin{tikzcd}[ampersand replacement=\&]
1\arrow[r]\&L_0\arrow[r]\&\widetilde B_0^*\arrow[r,"\det"]\&\unsure{\mathbf G_m|G_m}\arrow[r]\&1.
\end{tikzcd}}
\margin{il faudrait \unsure{vérifier} \illegible{4}\ \deleted{\illegible{1}}\quad\unsure{C'est} \uline{faux}, \unsure{cf plus bas}.}


\scanpage{276}
\mspage{439}

Soit $\underline u$ une section de $\widetilde B_0^*$, j'ai envie de prouver que si on pose
\[
(63)\qquad \xi=[\underline u,\rho],\quad\eta=[\underline u,\sigma],
\]
on aura automatiquement la relation
\[
(64)\qquad\xi=\eta\varepsilon_1(\mu-1)\qquad\text{où }\mu=\det\underline u.
\]
Pour nous \unsure{rassurer}, prenons d'abord $\underline u$ dans $L_0$, i.e.
\[
\underline u=\varepsilon_0(\lambda)
\]
d'où
\[
\begin{aligned}
\relax
[\underline u,\rho]&=\varepsilon_0(\lambda)\rho\varepsilon_0(\lambda)^{-1}\rho^{-1}
=\varepsilon_0(\lambda)\varepsilon_1(-\lambda)\\
[\underline u,\sigma]&=\varepsilon_0(\lambda)\sigma\varepsilon_0(-\lambda)\sigma^{-1}
\end{aligned}
\]
Or
\[
\sigma\underbrace{\varepsilon_0(-\lambda)}_{\inserted{1+N_0\lambda}}\sigma^{-1}
=1+(\sigma\uline{N_0}\sigma^{-1})\lambda\ \deleted{\unsure{=}}
\]
\unsure{or} \deleted{\illegible{1}} en vertu de $(57)$ on a
\[
\sigma N_0\sigma^{-1}=-(\lambda_0\rho\sigma^{-1}+1)\overset{(49)}{=}N_1
\]
donc on a
\[
(65)\quad\left\{\begin{array}{ll}
\sigma\varepsilon_0(\lambda)\sigma^{-1}=\varepsilon_1(\lambda)&\text{et pareillement}\\
\rho\varepsilon_0(\lambda)\rho^{-1}=\varepsilon_1(\lambda)
\end{array}\right.
\]
donc on trouve bien
\[
[\underline u,\rho]=[\underline u,\sigma]=\varepsilon_0(\lambda)\varepsilon_1(-\lambda)
\]
i.e. $(64)$, puisque ici $\unsure{\mu}=\det u=\unsure{1}$, donc $\mu-1=0$.


\scanpage{277}

Nous avons vu que $\rho$ et $\sigma$ tous les deux \unsure{transforment} $N_0$ dans $N_1$, donc le sous-groupe à $1$ \unsure{paramètre} additif
\[
L_0=\left\{\underset{\inserted{=1+\lambda N_0}}{\unsure{\varepsilon_0}(\lambda)}\ \middle|\ \lambda\text{ dans }\unsure{\mathbf G_a|G_a}\right\}
\]
dans le sous-groupe analogue
\[
L_1=\left\{\underset{\inserted{=1+\lambda N_1}}{\varepsilon_1(\lambda)}\ \middle|\ \lambda\text{ dans }\unsure{\mathbf G_a|G_a}\right\},
\]
et donc $\widetilde B_0$ dans $\widetilde B_1$ (défini en termes de $L_1$, ou de $N_1$, comme $\widetilde B_0^*$ en termes de $L_0$). \deleted{Donc} La formule $(\unsure{64|69})$ s'écrit aussi
\margin{\deleted{$\rho^{-1}\sigma$}}
\[
\rho\underline u^{-1}\rho^{-1}=\unsure{\sigma}\underline u^{-1}\sigma^{-1}\varepsilon_1(\mu-1)
\]
ou aussi, en remplaçant $\underline u$ par $\underline u^{-1}$
\[
(\unsure{*|**})\qquad
\rho\underline u\rho^{-1}=\sigma\underline u\sigma^{-1}\varepsilon_1\left(\frac1\mu-1\right).
\]
Or on \deleted{\illegible{1}} a
\[
\rho\underline u\rho^{-1}=\deleted{\unsure{\lambda}}\operatorname{int}(\lambda)(\sigma\underline u\sigma^{-1})
\]
où $\uline{\lambda=\rho\sigma^{-1}}$, et la formule $\unsure{(*)|(**)}$ s'écrit aussi
\[
(**):\qquad
\operatorname{int}(\lambda)(v)=v\cdot\varepsilon_1\left(\frac1\mu-1\right)
\qquad\begin{gathered}
\text{pour }v\ \uline{\text{dans }\unsure{B'_1|\widetilde B_1}}\\
\text{et }\mu=\det v;
\end{gathered}
\]
et en vertu de $(40\ \mathrm{bis})$ $\operatorname{int}\lambda=\operatorname{int}(\varepsilon_1)$, donc la formule se réduit à
\[
\operatorname{int}(\varepsilon_1)(v)=v\cdot\varepsilon_1\left(\frac1\mu-1\right)
\]


\scanpage{278}
\mspage{440}

ou encore
\[
\deleted{(1+N_1)\unsure{v_0|v}(1-N_1)=v\left(1+N_1\left(\frac1\mu-1\right)\right)}
\]
\deleted{soit}
\[
\deleted{v+N_1v-vN_1\unsure{+|-}N_1vN_1=v+vN_1\left(\frac1\mu-1\right)}
\]
\[
\varepsilon_1v\varepsilon_1^{\deleted{\unsure{-1}}}(-1)=v\varepsilon_1\left(\frac1\mu-1\right)\quad\text{soit}
\]
\[
\underbrace{\varepsilon_1}_{\inserted{1+N_1}}v
=v\underbrace{\varepsilon_1\left(\frac1\mu\right)}_{\inserted{1+\frac1\mu N_1}}
\qquad\text{soit }v+N_1v=v+\frac1\mu vN_1
\]
soit enfin
\[
vN_1v^{-1}=\mu N_1\quad(=\det(v).N_1)
\]
ce qui n'est autre que la définition de $\widetilde B_1^{\unsure{*}}$ !

On trouve donc un \deleted{\unsure{mono}}morphisme
\[
(66)\quad\left\{\begin{array}{rcl}
\widetilde B_0&\longrightarrow&\mathcal X_{\rho,\sigma}\\
\underline u&\longmapsto&([\underline u,\rho],[\underline u,\sigma],\det\underline u)
\end{array}\right.
\]
de $\widetilde B_0^{\unsure{*}}$ dans l'ens des solutions du système $(33)$ à $(35)$, est-ce un isomorphisme ?

\unsure{Notons} que l'égalité $[\underline u,\rho]=[\underline u',\rho]$ signifie que $u^{-1}\unsure{u'}\in\Gamma T_\rho$ donc \deleted{\illegible{3}} dire que $(66)$ est mono. signifie simplement que
\[
\deleted{(\unsure{67})}\qquad T_\rho\cap\widetilde B_0=\{1\}
\]
en fait l'analyse faite doit \unsure{montrer}

\scanpage{279}

que cette intersection est égale aussi à $T_\sigma\cap\widetilde B_0^{\unsure{*|\prime}}$, donc \unsure{contenue} dans $T_\sigma\cap T_\rho$. Je vais supposer
\[
\tag{68}
T_\sigma\cap T_\rho=\mathbf G_m\;(=\unsure{\mathrm{Cent}}\,G(M))
\]
\inserted{(\unsure{est} $\unsure{\supset|\supseteq}$, il \illegible{2} $\unsure{\subset|=}$)}
--- il \unsure{suffira|suffit} pour ceci de supposer qu'en chaque fibre, $1,\sigma,\rho$ sont linéairement indépendants --- c'est le cas bien sûr dans le cas qui nous intéresse. \uline{Mais} Donc
\[
T_\sigma\cap T_\rho\cap\widetilde B_0^{\unsure{*|\prime}}\subset\mathbf G_m\cap\widetilde B_0^*,
\]
mais
\[
\mathbf G_m\cap\widetilde B_0^*=\deleted{\{1\}}\mu_2
\]
Car si $\lambda$ \deleted{\illegible{1}} \unsure{est} une section \unsure{locale} de $\mathbf G_m$, alors
\[
\lambda N_0\lambda^{-1}=N_0=1.\unsure{N_0|N_3}
\]
mais $\det\lambda=\lambda^2$, et la condition que $\lambda$ \deleted{\unsure{soit}} (considéré comme section de $B_0$) soit dans $B_0^{\unsure{*|\prime}}$, c'est bien \uline{$\lambda^2=1$}. \deleted{\unsure{$B_0$}}

On trouve \deleted{\illegible{1}} ainsi
\[
\tag{69}
T_\sigma\cap\widetilde B_0^{\unsure{*|\prime}}=T_\rho\cap\widetilde B_0^{\unsure{*|\prime}}=\mu_2
\]
et \unsure{le}\deleted{\illegible{1}} morphisme (66) n'est pas un monomorphisme, mais se \uline{factorise}

\scanpage{280}
\mspage{441}

en un monomorphisme
\[
\tag{70}
\widetilde B_0/\mu_2\hookrightarrow\mathcal X_{\rho,\sigma}.
\]
À vrai dire, on voit que $\widetilde B_0$ n'est pas le ``bon'' groupe à introduire ici --- dans le cas particulier $\varepsilon_0=\begin{pmatrix}1&-1\\0&1\end{pmatrix}$, \deleted{\unsure{$B_0$}} on a
\[
B_0=\left\{\begin{pmatrix}p&\unsure{\ell|l}\\0&q\end{pmatrix}\mid pq\ \mathrm{inv.}\right\}
\]
et la relation \deleted{\unsure{pour}} $\chi'(\underline u)=\underbrace{\chi(\underline u)}_{\inserted{\det\underline u=pq}}$ s'écrit, par un calcul immédiat
\[
\tag{71}
p/q=pq\qquad\mathrm{i.e.}\quad q^2=1
\]
tandis que le ``bon'' groupe qui nous intéresse est défini par
\[
B'_0=\left\{\begin{pmatrix}p&\unsure{\ell|e}\\0&1\end{pmatrix}\mid p\ \mathrm{inv}\right\},
\]
\unsure{donc} avec $q=1$. Pour donner la description dans le cas général, on note que la fonction déterminant sur $B_0^*$ se décompose intrinsèquement en deux facteurs
\[
\tag{72}
\det\underline u=p(\underline u)q(\underline u),
\]
qu'on peut distinguer l'un de l'autre \unsure{justement} par la condition que
\[
\tag{73}
\underline uN_0\underline u^{-1}=\frac{p(\underline u)}{q(\underline u)}N_0,
\]
et la condition qui va nous décrire le ``bon'' groupe n'est pas (71), mais bien
\[
\tag{74}
q=q(\underline u)=1,
\]

\scanpage{281}

qui définit donc un groupe qui est \uline{lisse} ($\overline{\mathrm m}$ en car. 2 !) et s'insère cette fois-ci dans une suite exacte canonique
\[
\tag{74}
1\longrightarrow\unsure{L_0^*|L_0}\longrightarrow B'_0\xrightarrow{\det}\mathbf G_m\longrightarrow1
\]
On a maintenant
\[
\tag{75}
B'_0\cap\mathbf G_m=\{1\},
\]
ce qui implique que le morphisme induit par (66)
\[
\tag{76}
B'_0\hookrightarrow\mathcal X_{\rho,\sigma}
\]
est un monomorphisme.

Pour montrer que c'est un iso, il suffit maintenant, en vertu \unsure{des principes généraux}, de montrer que c'est \uline{surjectif} sur les pts à val. dans un corps alg clos (\unsure{si ça} \unsure{peut|pour} nous simplifier la vie). Je laisse tomber le détail \unsure{des|de} vérifications.

\deleted{\unsure{Rem}}

\uline{\unsure{Remq|Rem}} On peut considérer, par la condition (40 bis), que l'on part d'une section \inserted{$\varepsilon_1=1+N_1\quad(N_1^2=0)$} \deleted{\unsure{$\varepsilon_1$}} unipotente de $SG(M)$, et d'une section $\rho$, et que $\sigma$ en est déduit par \deleted{\unsure{une}}\inserted{la}

\scanpage{282}
\mspage{442}

formule
\[
\tag{77}
\sigma=-\varepsilon_1^{-1}\rho\;\left(=-(1-N_1)\rho\quad\mathrm{i.e.}\quad\sigma=-\rho+N_1\rho,\quad\mathrm{i.e.}\quad\rho+\sigma=N_1\rho\right)
\]

\margin{Je prends la constante qui \unsure{intervient} \unsure{dans|ds} le cas qui nous \unsure{intéresse} le plus, où $-\sigma^{-1}=\sigma$, de sorte que (40 bis) devient $\varepsilon_1=\rho\sigma$.}

(je laisse tomber la constante \inserted{\unsure{pure|pour}} $\lambda_0$, car \deleted{\illegible{1}} la situation n'est pas changée quand on remplace $\rho,\sigma$ par des multiples scalaires, car $F_\rho$ et $F_\sigma$ n'est pas changé \ldots) La condition (42) devient alors
\[
c=\operatorname{Tr}\rho+\operatorname{Tr}\sigma=\operatorname{Tr}(1-\underbrace{\varepsilon_1^{-1}}_{\inserted{1-N_1}})\rho=\operatorname{Tr}N_1\rho\in\Gamma\mathcal O_S^*
\]

\margin{\uline{NB} on a
\[
\tag{78 bis}
\left\{\begin{aligned}
\rho+\sigma&=N_1\rho=\rho N_0\\
&=-N_1\sigma=-\sigma N_0
\end{aligned}\right.
\]
\unsure{soit} $\unsure{\gamma_0|\varepsilon_0}$ \illegible{4}}

(section \uline{inversible} de $\mathcal O_S$). On suppose que $1,\rho,\sigma$ sont linéairement indépendants en chaque fibre, il revient au $\overline{\mathrm m}$ de dire que $1,\rho$, et $N_1\rho$ le sont, ou encore que $\rho^{-1},1$ et $N_1$ le sont, ou que $1,\sigma,N_1\sigma$ le sont, ou $\sigma^{-1},1,N_1$. En résumé

\margin{\unsure{NB} (Quitte à multiplier \unsure{le|les} $\rho$ par un scalaire inv) \unsure{on} \uline{$\operatorname{Tr}\unsure{\gamma_0|\delta_0}=1$} ($\unsure{\gamma_0|\delta_0}=N_1\rho$) ! (ce qui normalise \unsure{entièrement}, avec la relation $\varepsilon_1=-\rho\sigma^{-1}$, les facteurs scalaires \unsure{arbitraires|intrinsèques} de $\rho,\sigma$)}

\uline{Proposition}. \unsure{Soit} $\unsure{M|\mathcal M}$ une alg. d'Azumaya sur le schéma de base $S$, $N_1$ une section de carré nul \deleted{\unsure{qui}} i.e. $\varepsilon_1=1+N_1$ une section unipotente, $\rho$, \inserted{$\sigma$} \deleted{\unsure{une}} \inserted{des} sections reliées par
\[
\tag{\unsure{40|78}}
\varepsilon_1=-\rho\sigma^{-1}\quad\mathrm{i.e.}\quad\sigma=-\varepsilon_1^{-1}\rho\quad\mathrm{i.e.}\quad\rho+\sigma=N_1\rho\quad\mathrm{i.e.}\quad\unsure{\sigma+\rho|\rho+\sigma}=-N_1\sigma
\]
\deleted{\illegible{1}} (ce qui détermine donc chacune des trois sections \unsure{\deleted{$\varepsilon_1$}|$\varepsilon_1$} $\rho,\sigma$ et $\varepsilon_1$ (ou $N_1$) en fonction des deux autres, et on peut se donner arbitrairement $\rho$ et \deleted{\unsure{$N_1$}} $\varepsilon_1$ i.e. $\rho$ et $N_1$ de carré nul) --- on
